\section{Conclusão}

O desenvolvimento do sistema embarcado do Logan foi conduzido com foco na compatibilização entre requisitos regulatórios, necessidades operacionais das missões e limitações práticas de integração física e computacional da aeronave. Ao longo do projeto, buscou-se estruturar uma arquitetura capaz de integrar controle de voo, processamento embarcado, sensoriamento, comunicação e alimentação de forma coesa, confiável e aderente ao contexto da competição. 

As escolhas de componentes e a disposição física adotada não foram tratadas apenas como decisões de montagem, mas como definições de engenharia orientadas por critérios de desempenho, robustez, acessibilidade, compatibilidade eletromagnética e facilidade de manutenção. Nesse sentido, a arquitetura proposta procurou responder diretamente aos desafios identificados em ciclos anteriores, incorporando melhorias no arranjo dos módulos, na confiabilidade da navegação e na organização geral do sistema.

Por fim, o trabalho desenvolvido pelo setor de sistemas embarcados consolidou uma base técnica coerente para sustentar a operação do Logan nas missões da EletroQuad SAE Brasil 2026. A solução resultante reúne os elementos necessários para execução supervisionada das tarefas autônomas, preservando integração entre percepção, decisão, atuação e segurança operacional.
